\documentclass[11pt]{article}
\usepackage[margin=1in]{geometry}
\usepackage{amsmath,amssymb}
\usepackage{enumitem}
\usepackage{hyperref}
\usepackage{longtable}
\usepackage{array}
\usepackage{booktabs}
\usepackage{xcolor}
\usepackage{titlesec}
\usepackage{fancyhdr}

\titleformat{\section}{\large\bfseries}{\thesection}{1em}{}
\titleformat{\subsection}{\normalsize\bfseries}{\thesubsection}{1em}{}

\pagestyle{fancy}
\fancyhf{}
\rhead{Databricks ML Associate --- 100 Exam-Style Questions}
\lhead{Questions}
\cfoot{\thepage}

\title{\textbf{100 Exam-Style Practice Questions}\\[4pt]
	\large Databricks Certified \textbf{Machine Learning Associate}\\[2pt]
	\normalsize 4 options per question. Answer key with explanations on the final pages.}
\author{Practice Exam}
\date{}

\begin{document}
	\maketitle
	\tableofcontents
	\newpage
	
	\section*{Instructions}
	\addcontentsline{toc}{section}{Instructions}
	Each question has exactly four options (A--D). Choose the single best answer. The answer key with a short explanation of why the correct option is right and why the other three are wrong begins on page~\pageref{sec:key}.
	
	\textbf{Domains:} A. Databricks Machine Learning (Q1--38); B. ML Workflows (Q39--57); C. Model Development (Q58--88); D. Model Deployment (Q89--100).
	\newpage
	
	\section{A. Databricks Machine Learning}
	
	\begin{enumerate}[leftmargin=*, label=\textbf{Q\arabic*.}]
		
		\item A machine learning engineer created a Feature Table \texttt{new\_table} using Feature Store Client \texttt{fs} and specified a metadata description. Which line of code returns that description?
		\begin{enumerate}[label=(\alph*)]
			\item \texttt{fs.create\_training\_set("new\_table")}
			\item \texttt{fs.get\_table("new\_table").description}
			\item \texttt{fs.get\_table("new\_table").load\_df()}
			\item \texttt{fs.get\_table("new\_table")}
		\end{enumerate}
		
		\item What is the primary advantage of creating Feature Store tables at the account level in Unity Catalog rather than at the workspace level?
		\begin{enumerate}[label=(\alph*)]
			\item Faster model training
			\item Cross-workspace governance, discoverability, and access control
			\item Automatic feature engineering
			\item Lower storage cost
		\end{enumerate}
		
		\item Which client is used to create and manage Feature Store tables in Unity Catalog?
		\begin{enumerate}[label=(\alph*)]
			\item \texttt{FeatureStoreClient}
			\item \texttt{FeatureEngineeringClient}
			\item \texttt{UnityCatalogClient}
			\item \texttt{MLflowClient}
		\end{enumerate}
		
		\item Which method creates a feature table with \texttt{FeatureEngineeringClient}?
		\begin{enumerate}[label=(\alph*)]
			\item \texttt{fe.create\_table()}
			\item \texttt{fe.write\_table()}
			\item \texttt{fe.create\_training\_set()}
			\item \texttt{fe.log\_model()}
		\end{enumerate}
		
		\item Which method writes features to an existing feature table?
		\begin{enumerate}[label=(\alph*)]
			\item \texttt{fe.create\_table()}
			\item \texttt{fe.write\_table()}
			\item \texttt{fe.score\_batch()}
			\item \texttt{fe.get\_table()}
		\end{enumerate}
		
		\item What does \texttt{fe.create\_training\_set()} provide that a plain DataFrame join does not?
		\begin{enumerate}[label=(\alph*)]
			\item Automatic one-hot encoding
			\item Point-in-time lookups aligned to the label timestamp
			\item Model registration
			\item Hyperparameter tuning
		\end{enumerate}
		
		\item A model performs well offline but degrades in production. Review shows a feature was computed using data available only after the label outcome. What is the diagnosis?
		\begin{enumerate}[label=(\alph*)]
			\item Underfitting
			\item Training-serving skew / target leakage
			\item Class imbalance
			\item Concept drift
		\end{enumerate}
		
		\item What is the difference between online and offline feature tables?
		\begin{enumerate}[label=(\alph*)]
			\item Online is for training; offline is for serving
			\item Online serves low-latency real-time inference; offline is for batch training and scoring
			\item They are identical
			\item Offline is deprecated
		\end{enumerate}
		
		\item A feature-aware model was logged with \texttt{fs.log\_model()} and a \texttt{batch\_df} contains only the primary key \texttt{customer\_id}. Which call generates predictions by automatically joining the required features?
		\begin{enumerate}[label=(\alph*)]
			\item \texttt{model.predict(batch\_df)}
			\item \texttt{fs.score\_batch(model\_uri, batch\_df)}
			\item \texttt{mlflow.pyfunc.load\_model(model\_uri).predict(batch\_df)}
			\item \texttt{batch\_df.join(feature\_table)}
		\end{enumerate}
		
		\item Which namespace format is correct for a feature table in Unity Catalog?
		\begin{enumerate}[label=(\alph*)]
			\item \texttt{catalog.schema.table}
			\item \texttt{database.table}
			\item \texttt{workspace.table}
			\item \texttt{hive\_metastore.table}
		\end{enumerate}
		
		\item Which MLflow API call manually logs a metric?
		\begin{enumerate}[label=(\alph*)]
			\item \texttt{mlflow.log\_metric()}
			\item \texttt{mlflow.log\_param()}
			\item \texttt{mlflow.log\_artifact()}
			\item \texttt{mlflow.set\_tag()}
		\end{enumerate}
		
		\item Which MLflow API call registers a model in Unity Catalog?
		\begin{enumerate}[label=(\alph*)]
			\item \texttt{mlflow.log\_model()}
			\item \texttt{mlflow.register\_model()}
			\item \texttt{mlflow.save\_model()}
			\item \texttt{mlflow.start\_run()}
		\end{enumerate}
		
		\item What replaced the old Staging and Production stages for model promotion in Unity Catalog?
		\begin{enumerate}[label=(\alph*)]
			\item Tags only
			\item Aliases such as \texttt{@champion}
			\item Notebook revisions
			\item Git branches
		\end{enumerate}
		
		\item How do you promote a challenger model to champion using aliases?
		\begin{enumerate}[label=(\alph*)]
			\item Delete the old model version
			\item Move the \texttt{champion} alias to the new model version
			\item Rename the model
			\item Create a new experiment
		\end{enumerate}
		
		\item What does \texttt{mlflow.search\_runs()} return?
		\begin{enumerate}[label=(\alph*)]
			\item A single run ID
			\item A pandas DataFrame of runs that can be sorted and filtered
			\item A trained model
			\item A list of experiments
		\end{enumerate}
		
		\item A run ID is stored in \texttt{run\_id} and the logged model name is \texttt{"model"}. Which line registers it as \texttt{"best\_model"}?
		\begin{enumerate}[label=(\alph*)]
			\item \texttt{mlflow.register\_model(run\_id, "best\_model")}
			\item \texttt{mlflow.register\_model(f"runs:/\{run\_id\}/model", "best\_model")}
			\item \texttt{mlflow.register\_model(f"runs:/\{run\_id\}/model")}
			\item \texttt{mlflow.register\_model(f"runs:/\{run\_id\}/best\_model", "model")}
		\end{enumerate}
		
		\item A data scientist wants one parent run for tuning and a child run per trial, starting runs manually. What should they do?
		\begin{enumerate}[label=(\alph*)]
			\item Turn on Databricks Autologging
			\item Specify \texttt{nested=True} when starting each child run
			\item Start each child run inside the parent's code block without \texttt{nested}
			\item Give child runs the same experiment ID
		\end{enumerate}
		
		\item Which UI element on the MLflow experiment page opens the notebook that created a run?
		\begin{enumerate}[label=(\alph*)]
			\item The \texttt{MLmodel} artifact
			\item The \texttt{Models} link
			\item The \texttt{Source} link
			\item The \texttt{Start Time} link
		\end{enumerate}
		
		\item Which function searches runs ordered by a metric ascending to find the lowest RMSE?
		\begin{enumerate}[label=(\alph*)]
			\item \texttt{mlflow.search\_runs(experiment\_id, order\_by=["metrics.rmse DESC"])}
			\item \texttt{mlflow.best\_run(experiment\_id, order\_by=["metrics.rmse"])}
			\item \texttt{mlflow.search\_runs(experiment\_id, order\_by=["metrics.rmse ASC"])}
			\item \texttt{mlflow.search\_runs(experiment\_id, filter="rmse < 1")}
		\end{enumerate}
		
		\item On which page do you transition a registered model version from Staging to Production?
		\begin{enumerate}[label=(\alph*)]
			\item The MLflow Model Registry home page
			\item The experiment page
			\item The model version page
			\item The model page only
		\end{enumerate}
		
		\item Which MLflow API call logs an artifact?
		\begin{enumerate}[label=(\alph*)]
			\item \texttt{mlflow.log\_metric()}
			\item \texttt{mlflow.log\_param()}
			\item \texttt{mlflow.log\_artifact()}
			\item \texttt{mlflow.set\_tag()}
		\end{enumerate}
		
		\item What is the simplest Databricks-native way to make every training run reproducible months later?
		\begin{enumerate}[label=(\alph*)]
			\item Save notebooks to a local drive
			\item Use MLflow tracking to log parameters, metrics, artifacts, and the model
			\item Print metrics to stdout
			\item Rely on cluster logs
		\end{enumerate}
		
		\item What does Databricks AutoML produce besides a trained model?
		\begin{enumerate}[label=(\alph*)]
			\item Only a model artifact
			\item A leaderboard, editable notebooks, and an automated EDA notebook
			\item A production endpoint
			\item A data warehouse
		\end{enumerate}
		
		\item Which API runs classification AutoML?
		\begin{enumerate}[label=(\alph*)]
			\item \texttt{databricks.automl.classify()}
			\item \texttt{databricks.automl.regress()}
			\item \texttt{databricks.automl.forecast()}
			\item \texttt{databricks.automl.train()}
		\end{enumerate}
		
		\item AutoML is least appropriate when:
		\begin{enumerate}[label=(\alph*)]
			\item You need a quick baseline
			\item You need a highly customized deep learning architecture
			\item You want automatic feature engineering
			\item You want to compare multiple algorithms
		\end{enumerate}
		
		\item Which step must a data scientist perform outside an AutoML experiment?
		\begin{enumerate}[label=(\alph*)]
			\item Model tuning
			\item Model evaluation
			\item Model deployment
			\item Exploratory data analysis
		\end{enumerate}
		
		\item How do you begin using Databricks Runtime ML?
		\begin{enumerate}[label=(\alph*)]
			\item Add an init script line
			\item Check a Databricks Runtime ML box
			\item Select a Runtime ML version from the Databricks Runtime Version dropdown
			\item Set \texttt{runtime-version="ml"} in the Spark session
		\end{enumerate}
		
		\item Why is Databricks Runtime ML convenient for MLflow users?
		\begin{enumerate}[label=(\alph*)]
			\item It removes the need for a cluster
			\item It preinstalls MLflow and other ML libraries
			\item It disables logging
			\item It only supports R
		\end{enumerate}
		
		\item Which evaluation metric is not suitable for AutoML regression experiments?
		\begin{enumerate}[label=(\alph*)]
			\item F1
			\item R-squared
			\item MAE
			\item MSE
		\end{enumerate}
		
		\item Which file typically defines a model's format and dependencies in an MLflow run?
		\begin{enumerate}[label=(\alph*)]
			\item \texttt{requirements.txt}
			\item \texttt{MLmodel}
			\item \texttt{README.md}
			\item \texttt{conda.yaml} only
		\end{enumerate}
		
		\item What does \texttt{fs.log\_model()} add to a model artifact?
		\begin{enumerate}[label=(\alph*)]
			\item Hyperparameters
			\item Feature lookup metadata for automatic joins at inference
			\item A confusion matrix
			\item A serving endpoint
		\end{enumerate}
		
		\item Which is a benefit of registering models in Unity Catalog over the workspace registry?
		\begin{enumerate}[label=(\alph*)]
			\item Faster training
			\item Lineage, governance, and cross-workspace access
			\item Automatic hyperparameter tuning
			\item Free storage
		\end{enumerate}
		
		\item What does an MLflow experiment group together?
		\begin{enumerate}[label=(\alph*)]
			\item Runs of a related ML task
			\item Clusters
			\item Notebooks only
			\item Feature tables
		\end{enumerate}
		
		\item Which API call starts an MLflow run?
		\begin{enumerate}[label=(\alph*)]
			\item \texttt{mlflow.start\_run()}
			\item \texttt{mlflow.log\_run()}
			\item \texttt{mlflow.open\_run()}
			\item \texttt{mlflow.new\_run()}
		\end{enumerate}
		
		\item Which is true about aliases in the Unity Catalog model registry?
		\begin{enumerate}[label=(\alph*)]
			\item An alias can point to many versions at once
			\item An alias points to one model version and can be moved
			\item Aliases are immutable
			\item Aliases replace experiments
		\end{enumerate}
		
		\item Which object stores the actual model files and metadata for a run?
		\begin{enumerate}[label=(\alph*)]
			\item The artifact store
			\item The parameter log
			\item The metric log
			\item The experiment tag
		\end{enumerate}
		
		\item Why should secrets not be logged as MLflow parameters?
		\begin{enumerate}[label=(\alph*)]
			\item Parameters are stored in plain text and visible to experiment users
			\item Parameters are not persisted
			\item Parameters are encrypted by default
			\item Parameters are deleted after the run
		\end{enumerate}
		
	\end{enumerate}
	
	\newpage
	\section{B. ML Workflows}
	
	\begin{enumerate}[leftmargin=*, label=\textbf{Q\arabic*.}, start=39]
		
		\item A Spark DataFrame \texttt{spark\_df} must keep only rows where \texttt{price > 0}. Which code block works?
		\begin{enumerate}[label=(\alph*)]
			\item \texttt{spark\_df[spark\_df["price"] > 0]}
			\item \texttt{spark\_df.filter(col("price") > 0)}
			\item \texttt{SELECT * FROM spark\_df WHERE price > 0}
			\item \texttt{spark\_df.loc[spark\_df["price"] > 0,:]}
		\end{enumerate}
		
		\item Which line returns count, mean, standard deviation, min, max, and quantiles for numeric columns?
		\begin{enumerate}[label=(\alph*)]
			\item \texttt{spark\_df.summary()}
			\item \texttt{spark\_df.stats()}
			\item \texttt{spark\_df.describe().head()}
			\item \texttt{spark\_df.printSchema()}
		\end{enumerate}
		
		\item How do you remove outliers using IQR in Spark?
		\begin{enumerate}[label=(\alph*)]
			\item Drop the top 10\% of rows
			\item Filter values outside $[Q1 - 1.5 \cdot IQR,\; Q3 + 1.5 \cdot IQR]$
			\item Filter values outside $[Q1 - IQR,\; Q3 + IQR]$
			\item Replace all values with the mean
		\end{enumerate}
		
		\item When is median imputation preferred over mean imputation?
		\begin{enumerate}[label=(\alph*)]
			\item When data is symmetric
			\item When data is skewed or contains extreme outliers
			\item When there are no missing values
			\item Always
		\end{enumerate}
		
		\item A data scientist replaces missing values with the median. A colleague says valuable information is lost. Which approach preserves the most information?
		\begin{enumerate}[label=(\alph*)]
			\item Impute with the mean instead
			\item Let the algorithm handle missing values
			\item Drop all features with missing values
			\item Add a binary indicator feature marking which values were imputed
		\end{enumerate}
		
		\item One-hot encoding is least appropriate for:
		\begin{enumerate}[label=(\alph*)]
			\item Nominal categorical features
			\item Very high-cardinality categorical features
			\item Low-cardinality categorical features
			\item Tree-based models
		\end{enumerate}
		
		\item A feature repository one-hot encodes all categorical variables. Why is this problematic?
		\begin{enumerate}[label=(\alph*)]
			\item One-hot encoding is not supported by most libraries
			\item It depends on target values
			\item It is computationally intensive and should be done on small samples
			\item It is a potentially problematic strategy for some ML algorithms due to dimensionality and sparsity
		\end{enumerate}
		
		\item How do you one-hot encode in Spark ML?
		\begin{enumerate}[label=(\alph*)]
			\item \texttt{OneHotEncoder} alone
			\item \texttt{StringIndexer} then \texttt{OneHotEncoder}
			\item \texttt{VectorAssembler} then \texttt{StringIndexer}
			\item \texttt{Imputer} then \texttt{OneHotEncoder}
		\end{enumerate}
		
		\item What does \texttt{VectorAssembler} do?
		\begin{enumerate}[label=(\alph*)]
			\item Combines multiple columns into a single feature vector
			\item Splits a vector into columns
			\item Encodes strings
			\item Imputes missing values
		\end{enumerate}
		
		\item A data scientist creates an \texttt{Imputer} and calls \texttt{transform} directly on it. Why does it fail?
		\begin{enumerate}[label=(\alph*)]
			\item It does not impute both training and test sets
			\item \texttt{inputCols} and \texttt{outputCols} must be identical
			\item It should call \texttt{fit} instead of \texttt{transform}
			\item It does not fit the imputer to create an \texttt{ImputerModel}
		\end{enumerate}
		
		\item When is log transformation appropriate?
		\begin{enumerate}[label=(\alph*)]
			\item For right-skewed positive features
			\item For left-skewed negative features
			\item For categorical features
			\item For binary features
		\end{enumerate}
		
		\item A pandas notebook must scale to big data with the least refactoring effort. What approach is recommended?
		\begin{enumerate}[label=(\alph*)]
			\item Refactor to process data in parallel
			\item Refactor to the PySpark DataFrame API
			\item Refactor to the Scala Dataset API
			\item Refactor to Spark SQL
			\item Refactor to the pandas API on Spark
		\end{enumerate}
		
		\item Which function applies a Python function to each group of a grouped DataFrame and returns a DataFrame?
		\begin{enumerate}[label=(\alph*)]
			\item \texttt{applyInPandas}
			\item \texttt{groupedApplyInPandas}
			\item \texttt{mapInPandas}
			\item \texttt{predict}
		\end{enumerate}
		
		\item A data scientist wants group-specific model inference in parallel. Which call fills the blank in \texttt{df.groupby("device\_id").\_\_\_(apply\_model, schema=...) }?
		\begin{enumerate}[label=(\alph*)]
			\item \texttt{applyInPandas}
			\item \texttt{groupedApplyInPandas}
			\item \texttt{mapInPandas}
			\item \texttt{predict}
		\end{enumerate}
		
		\item Which imputation strategy is most appropriate for a numeric feature with extreme outliers?
		\begin{enumerate}[label=(\alph*)]
			\item Mean
			\item Median
			\item Mode
			\item Zero
		\end{enumerate}
		
		\item What does a boxplot show?
		\begin{enumerate}[label=(\alph*)]
			\item Only the mean
			\item Median, quartiles, and outliers
			\item Correlation
			\item Class balance
		\end{enumerate}
		
		\item Which technique helps detect multicollinearity?
		\begin{enumerate}[label=(\alph*)]
			\item Variance Inflation Factor
			\item Silhouette score
			\item ROC AUC
			\item Confusion matrix
		\end{enumerate}
		
		\item Why is stratification important for classification splits?
		\begin{enumerate}[label=(\alph*)]
			\item It preserves class proportions in each split
			\item It increases training speed
			\item It removes outliers
			\item It ensures reproducibility only
		\end{enumerate}
		
		\item Which is a symptom of an imbalanced dataset?
		\begin{enumerate}[label=(\alph*)]
			\item A 50/50 class split
			\item A 99\% negative / 1\% positive split
			\item A continuous target with normal distribution
			\item Missing values
		\end{enumerate}
		
	\end{enumerate}
	
	\newpage
	\section{C. Model Development}
	
	\begin{enumerate}[leftmargin=*, label=\textbf{Q\arabic*.}, start=58]
		
		\item Which algorithm typically uses bagging?
		\begin{enumerate}[label=(\alph*)]
			\item Gradient boosted trees
			\item K-means
			\item Random forest
			\item Linear regression
		\end{enumerate}
		
		\item Spark ML linear regression on large data uses which approach for distributed training?
		\begin{enumerate}[label=(\alph*)]
			\item Logistic regression
			\item Spark ML cannot distribute linear regression
			\item Iterative optimization such as L-BFGS
			\item Least-squares normal equations
		\end{enumerate}
		
		\item Why can a Spark ML decision tree differ from a scikit-learn tree with identical data and hyperparameters?
		\begin{enumerate}[label=(\alph*)]
			\item Spark ML tests every feature
			\item Spark ML automatically prunes overfit trees
			\item Spark ML tests more split candidates
			\item Spark ML tests a random sample of features
			\item Spark ML tests binned feature values as representative split candidates
		\end{enumerate}
		
		\item What does Hyperopt's \texttt{fmin} do?
		\begin{enumerate}[label=(\alph*)]
			\item Minimizes an objective function over a search space
			\item Maximizes accuracy only
			\item Trains a model
			\item Logs a metric
		\end{enumerate}
		
		\item What does \texttt{SparkTrials} do?
		\begin{enumerate}[label=(\alph*)]
			\item Distributes Hyperopt trials across a Spark cluster
			\item Runs trials on the driver only
			\item Saves models to the registry
			\item Encodes categorical features
		\end{enumerate}
		
		\item What is the role of \texttt{tpe.suggest} in Hyperopt?
		\begin{enumerate}[label=(\alph*)]
			\item Evaluates the objective function
			\item Specifies the MLflow experiment
			\item Suggests new hyperparameter configurations based on previous results
			\item Parallelizes the tuning process
		\end{enumerate}
		
		\item A team has 8 workers and 5 hyperparameters. Which approach best leverages the cluster?
		\begin{enumerate}[label=(\alph*)]
			\item Grid search on the driver
			\item Sequential manual tuning
			\item Hyperopt with \texttt{SparkTrials} and TPE
			\item Random search on the driver
		\end{enumerate}
		
		\item With 3-fold CV and a grid of 3 $\times$ 2 hyperparameter values, how many models are trained in total?
		\begin{enumerate}[label=(\alph*)]
			\item 6
			\item 9
			\item 18
			\item 36
		\end{enumerate}
		
		\item With 3-fold CV and a grid of 3 $\times$ 2 values, how many models can be trained in parallel?
		\begin{enumerate}[label=(\alph*)]
			\item 3
			\item 5
			\item 6
			\item 18
		\end{enumerate}
		
		\item What is a primary advantage of Hyperopt over manual tuning or grid search?
		\begin{enumerate}[label=(\alph*)]
			\item It efficiently explores the search space using algorithms like TPE
			\item Manual tuning is faster
			\item Grid search is more reliable
			\item Hyperopt needs more evaluations
		\end{enumerate}
		
		\item When is \texttt{StringIndexer} used in a pipeline?
		\begin{enumerate}[label=(\alph*)]
			\item To convert textual or categorical values to numeric indices
			\item To impute missing values
			\item To scale numeric features
			\item To assemble vectors
		\end{enumerate}
		
		\item Why might a train-validation split be more efficient than k-fold CV?
		\begin{enumerate}[label=(\alph*)]
			\item It requires testing fewer models, saving compute
			\item It reduces hyperparameter values to test
			\item It guarantees reproducibility
			\item It ensures bias-free training
		\end{enumerate}
		
		\item When using a pipeline and a cross-validator, where should they be placed?
		\begin{enumerate}[label=(\alph*)]
			\item Always put the cross-validator inside the pipeline
			\item If the pipeline includes estimators, put the entire pipeline inside the cross-validator
			\item If concerned about leakage, put the pipeline inside the cross-validator
			\item Always put the pipeline inside the cross-validator
		\end{enumerate}
		
		\item A health organization wants to maximize the number of positive infection cases identified. Which metric should be used?
		\begin{enumerate}[label=(\alph*)]
			\item RMSE
			\item Precision
			\item Area under the residual operating curve
			\item Accuracy
			\item Recall
		\end{enumerate}
		
		\item When should F1 be used over accuracy?
		\begin{enumerate}[label=(\alph*)]
			\item When classes are balanced
			\item When there are more than two classes
			\item When there is significant imbalance and avoiding false negatives is a priority
			\item When true positives and true negatives are equally important
		\end{enumerate}
		
		\item A fraud dataset has a 1\% positive class and a model reports 99\% accuracy. What should you do?
		\begin{enumerate}[label=(\alph*)]
			\item Deploy it
			\item Evaluate precision, recall, F1, and PR-AUC
			\item Increase epochs
			\item Add features
		\end{enumerate}
		
		\item What does $R^2 = 0.85$ mean?
		\begin{enumerate}[label=(\alph*)]
			\item 85\% of predictions are correct
			\item The model explains 85\% of the variance in the target
			\item The model has 85\% precision
			\item The RMSE is 0.85
		\end{enumerate}
		
		\item Why might RMSE be higher than MAE?
		\begin{enumerate}[label=(\alph*)]
			\item RMSE penalizes large errors more heavily
			\item MAE squares errors
			\item RMSE ignores outliers
			\item MAE is in different units
		\end{enumerate}
		
		\item Training accuracy is 99.5\% and validation accuracy is 72\%. What is the diagnosis and fix?
		\begin{enumerate}[label=(\alph*)]
			\item Underfitting; increase depth
			\item Overfitting; reduce complexity or add regularization
			\item Good model; deploy
			\item Need more test data
		\end{enumerate}
		
		\item What is the effect of increasing model complexity on bias and variance?
		\begin{enumerate}[label=(\alph*)]
			\item Bias decreases, variance increases
			\item Bias increases, variance decreases
			\item Both decrease
			\item Both increase
		\end{enumerate}
		
		\item A log(price) model is evaluated against actual price values and its RMSE is much larger than a price model. Which explanation is invalid?
		\begin{enumerate}[label=(\alph*)]
			\item The second model is much more accurate
			\item They failed to exponentiate predictions before computing RMSE
			\item They failed to take the log of predictions in the first model
			\item The first model is much more accurate
			\item RMSE is invalid for regression
		\end{enumerate}
		
		\item What is the 3-fold CV RMSE given fold RMSEs 10.0, 12.0, and 17.0?
		\begin{enumerate}[label=(\alph*)]
			\item 13.0
			\item 17.0
			\item 12.0
			\item 39.0
		\end{enumerate}
		
		\item What happens if the model is placed as a stage inside a pipeline and the cross-validator is the final stage?
		\begin{enumerate}[label=(\alph*)]
			\item Training takes longer
			\item Feature engineering is computed using validation data, causing leakage
			\item CV is no longer reproducible
			\item The model is refit once more per fold
		\end{enumerate}
		
		\item What does \texttt{CrossValidator} do in Spark ML?
		\begin{enumerate}[label=(\alph*)]
			\item Performs k-fold CV with a \texttt{ParamGridBuilder}
			\item Trains one model
			\item Clusters data
			\item Imputes missing values
		\end{enumerate}
		
		\item When is a train-validation split preferable to k-fold CV?
		\begin{enumerate}[label=(\alph*)]
			\item When high interpretability is required
			\item When compute is limited, especially for large data or complex models
			\item When the algorithm is highly parallelizable
			\item When compute is unlimited
		\end{enumerate}
		
		\item Which technique mitigates overfitting?
		\begin{enumerate}[label=(\alph*)]
			\item Feature engineering and dimensionality reduction
			\item Increasing model complexity
			\item Using only one technique such as data augmentation
			\item All of the above
		\end{enumerate}
		
		\item What is k-Nearest Neighbors' key hyperparameter tradeoff?
		\begin{enumerate}[label=(\alph*)]
			\item Small $k$ gives low bias and high variance; large $k$ gives high bias and low variance
			\item Small $k$ gives high bias; large $k$ gives high variance
			\item $k$ has no effect
			\item $k$ only affects training time
		\end{enumerate}
		
		\item Which kernel is commonly used for non-linear SVM classification?
		\begin{enumerate}[label=(\alph*)]
			\item Linear
			\item RBF
			\item Identity
			\item Zero
		\end{enumerate}
		
		\item What is gradient boosting?
		\begin{enumerate}[label=(\alph*)]
			\item Trees trained in parallel on bootstrap samples
			\item Trees trained sequentially, each correcting previous errors
			\item A clustering method
			\item A dimensionality reduction method
		\end{enumerate}
		
		\item What is the role of \texttt{max\_depth} in tree models?
		\begin{enumerate}[label=(\alph*)]
			\item Limits tree depth to control complexity and overfitting
			\item Sets the number of trees
			\item Sets the learning rate
			\item Sets the number of features
		\end{enumerate}
		
	\end{enumerate}
	
	\newpage
	\section{D. Model Deployment}
	
	\begin{enumerate}[leftmargin=*, label=\textbf{Q\arabic*.}, start=89]
		
		\item What are the three model serving approaches tested on the exam?
		\begin{enumerate}[label=(\alph*)]
			\item Batch, real-time, and streaming
			\item Batch, offline, and online
			\item Training, validation, and testing
			\item SQL, Python, and Scala
		\end{enumerate}
		
		\item A scikit-learn model must score 100 million Delta records efficiently. What is the best approach?
		\begin{enumerate}[label=(\alph*)]
			\item Load into pandas on the driver and loop
			\item Use \texttt{mlflow.pyfunc.spark\_udf()} to distribute inference
			\item Export to CSV and score externally
			\item Create a real-time endpoint and call per record
		\end{enumerate}
		
		\item How is streaming inference performed in Databricks?
		\begin{enumerate}[label=(\alph*)]
			\item Using Delta Live Tables
			\item Using a SQL Warehouse
			\item Using the driver only
			\item Using a pandas UDF on a single node
		\end{enumerate}
		
		\item How do you split traffic between real-time serving endpoints?
		\begin{enumerate}[label=(\alph*)]
			\item Configure traffic splitting on the serving endpoint
			\item Duplicate the model
			\item Use a SQL query
			\item Rename the endpoint
		\end{enumerate}
		
		\item What is the primary purpose of a real-time serving endpoint?
		\begin{enumerate}[label=(\alph*)]
			\item Batch scoring
			\item Low-latency predictions accessible via REST API
			\item Training models
			\item Data storage
		\end{enumerate}
		
		\item A team replaces one model with three new models of nearly equal latency. When is the new solution less efficient during inference?
		\begin{enumerate}[label=(\alph*)]
			\item When if-else logic picks a model per prediction
			\item When average model latency exceeds the original model's size
			\item When fewer features are used
			\item When every record must be scored by all three models
		\end{enumerate}
		
		\item What is a batch inference job in Databricks?
		\begin{enumerate}[label=(\alph*)]
			\item Scoring a large dataset offline and writing predictions to a table
			\item A real-time REST endpoint
			\item A training job
			\item A cluster configuration
		\end{enumerate}
		
		\item Which compute type is best for Spark SQL data prep plus Spark ML training in a notebook?
		\begin{enumerate}[label=(\alph*)]
			\item Single Node cluster
			\item Standard (all-purpose) cluster
			\item SQL Warehouse
			\item None of these
		\end{enumerate}
		
		\item Why is model monitoring important in production?
		\begin{enumerate}[label=(\alph*)]
			\item Data distributions and relationships can drift over time
			\item Once deployed, models never degrade
			\item It replaces evaluation
			\item It improves training speed
		\end{enumerate}
		
		\item What is data drift versus concept drift?
		\begin{enumerate}[label=(\alph*)]
			\item Data drift changes input distributions; concept drift changes the feature-target relationship
			\item Both refer to the same thing
			\item Data drift changes labels; concept drift changes features
			\item Neither affects models
		\end{enumerate}
		
		\item Which is a fairness concern when deploying ML?
		\begin{enumerate}[label=(\alph*)]
			\item Biased training data leading to discriminatory outcomes
			\item Faster inference
			\item Lower storage cost
			\item Better dashboards
		\end{enumerate}
		
		\item Which practice best guards against leakage in a production pipeline?
		\begin{enumerate}[label=(\alph*)]
			\item Fit all transforms on the entire dataset
			\item Encapsulate preprocessing and model in a Pipeline fit only on training data
			\item Use the test set for feature selection
			\item Scale before splitting
		\end{enumerate}
		
	\end{enumerate}
	
	\newpage
	\section{Answer Key with Explanations}
	\label{sec:key}
	
	\begin{enumerate}[leftmargin=*, label=\textbf{Q\arabic*.}]
		
		\item \textbf{B}. \texttt{get\_table()} returns the table object; \texttt{.description} reads metadata. (A) creates a training set; (C) loads data; (D) returns the object without the description.
		
		\item \textbf{B}. Account-level tables are governed and discoverable across workspaces. (A) and (D) are not the benefit; (C) is false.
		
		\item \textbf{B}. \texttt{FeatureEngineeringClient} is the Unity Catalog Feature Store client. (A) is the legacy client; (C) and (D) are wrong.
		
		\item \textbf{A}. \texttt{create\_table()} creates. (B) writes; (C) builds a training set; (D) logs a model.
		
		\item \textbf{B}. \texttt{write\_table()} writes features. (A) creates; (C) scores; (D) retrieves.
		
		\item \textbf{B}. Training sets use point-in-time lookups aligned to label timestamps. (A), (C), (D) are unrelated.
		
		\item \textbf{B}. Leakage from post-label data causes training-serving skew. (A), (C), (D) do not match the symptom.
		
		\item \textbf{B}. Online tables serve low-latency inference; offline tables are for batch. (A) reverses them; (C), (D) false.
		
		\item \textbf{B}. \texttt{fs.score\_batch()} uses logged feature metadata to join features automatically. (A), (C) assume features are already in the DataFrame; (D) is a plain join without automatic lookup.
		
		\item \textbf{A}. Unity Catalog uses \texttt{catalog.schema.table}. (B), (C), (D) are not Unity Catalog formats.
		
		\item \textbf{A}. \texttt{log\_metric()} logs metrics. (B) logs params; (C) artifacts; (D) tags.
		
		\item \textbf{B}. \texttt{register\_model()} registers. (A) logs; (C) saves; (D) starts a run.
		
		\item \textbf{B}. Aliases replace stages. (A), (C), (D) are not the replacement mechanism.
		
		\item \textbf{B}. Promotion is moving the alias. (A) deletes history; (C), (D) do not promote.
		
		\item \textbf{B}. It returns a pandas DataFrame of runs. (A) is too narrow; (C), (D) are false.
		
		\item \textbf{B}. The URI must include \texttt{runs:/} and the artifact path. (A) passes only the run ID; (C) omits the registry name; (D) swaps artifact path and name.
		
		\item \textbf{B}. Child runs must be started with \texttt{nested=True}. (A) does not nest; (C) raises an error; (D) only groups runs in the same experiment.
		
		\item \textbf{C}. The Source link opens the notebook snapshot. (A) is a model metadata file; (B) links to registry; (D) opens run details.
		
		\item \textbf{C}. \texttt{search\_runs} with ascending order returns the lowest RMSE first. (A) orders descending; (B) is not the API; (D) filters rather than sorts.
		
		\item \textbf{C}. Stage transitions happen on the model version page. (A), (B), (D) are not the direct page.
		
		\item \textbf{C}. \texttt{log\_artifact()} logs files. (A), (B), (D) log other things.
		
		\item \textbf{B}. MLflow tracking logs params, metrics, artifacts, and models. (A), (C), (D) are not reproducible logs.
		
		\item \textbf{B}. AutoML returns a leaderboard, notebooks, and an EDA notebook. (A) is incomplete; (C), (D) are not produced.
		
		\item \textbf{A}. Classification AutoML uses \texttt{classify()}. (B) is regression; (C) forecasting; (D) not an API.
		
		\item \textbf{B}. AutoML is less suited to highly custom deep learning. (A), (C), (D) are strengths.
		
		\item \textbf{C}. Deployment is outside AutoML. (A), (B), (D) are automated.
		
		\item \textbf{C}. Runtime ML is selected from the runtime dropdown. (A), (B), (D) do not select the runtime.
		
		\item \textbf{B}. Runtime ML preinstalls MLflow and ML libraries. (A), (C), (D) are false.
		
		\item \textbf{A}. F1 is for classification, not regression. (B), (C), (D) are regression metrics.
		
		\item \textbf{B}. \texttt{MLmodel} defines format and dependencies. (A), (D) are dependency specs; (C) is documentation.
		
		\item \textbf{B}. \texttt{log\_model()} packages feature lookup metadata. (A), (C), (D) are not added by this call.
		
		\item \textbf{B}. Unity Catalog adds lineage, governance, and cross-workspace access. (A), (C), (D) are false.
		
		\item \textbf{A}. Experiments group related runs. (B), (C), (D) are not grouped by experiments.
		
		\item \textbf{A}. \texttt{start\_run()} starts a run. (B), (C), (D) are not APIs.
		
		\item \textbf{B}. An alias points to one version and can be moved. (A) is false; (C) aliases are movable; (D) false.
		
		\item \textbf{A}. Artifacts live in the artifact store. (B), (C), (D) store metadata, not files.
		
		\item \textbf{A}. Parameters are stored in plain text and visible. (B), (C), (D) are false.
		
		\item \textbf{B}. \texttt{filter(col(...))} is the PySpark way. (A) is pandas-like; (C) is SQL, not DataFrame code; (D) is pandas.
		
		\item \textbf{A}. \texttt{summary()} returns count, mean, std, min, max, and quantiles. (B) is not an API; (C) lacks quantiles; (D) prints schema.
		
		\item \textbf{B}. The 1.5$\times$IQR rule is standard. (A), (C), (D) are wrong thresholds or methods.
		
		\item \textbf{B}. Median is robust to skew and outliers. (A), (C), (D) are wrong contexts.
		
		\item \textbf{D}. Indicator features preserve missingness information. (A), (B), (C) lose or ignore information.
		
		\item \textbf{B}. High cardinality causes dimensionality explosion. (A), (C), (D) are not the main problem.
		
		\item \textbf{D}. Premature one-hot can hurt some algorithms. (A), (B), (C) are false or irrelevant.
		
		\item \textbf{B}. Spark ML uses StringIndexer then OneHotEncoder. (A), (C), (D) are wrong order or wrong tools.
		
		\item \textbf{A}. VectorAssembler combines columns into a vector. (B), (C), (D) are different operations.
		
		\item \textbf{D}. The imputer must be fit first. (A), (B), (C) are not the cause.
		
		\item \textbf{A}. Log transform suits right-skewed positive data. (B), (C), (D) are wrong distributions.
		
		\item \textbf{E}. pandas API on Spark minimizes refactoring. (A), (B), (C), (D) require larger rewrites.
		
		\item \textbf{A}. \texttt{applyInPandas} applies per group and returns a DataFrame. (B) is not an API; (C) maps partitions; (D) is not an API.
		
		\item \textbf{A}. Same as above. (B), (C), (D) are wrong.
		
		\item \textbf{B}. Median is robust to outliers. (A), (C), (D) are not robust.
		
		\item \textbf{B}. Boxplots show median, quartiles, and outliers. (A), (C), (D) are not boxplot content.
		
		\item \textbf{A}. VIF detects multicollinearity. (B), (C), (D) measure other things.
		
		\item \textbf{A}. Stratification preserves class proportions. (B), (C), (D) are unrelated.
		
		\item \textbf{B}. 99/1 is severe imbalance. (A), (C), (D) are not imbalance symptoms.
		
		\item \textbf{C}. Random forest uses bagging. (A) uses boosting; (B) is clustering; (D) is a single model.
		
		\item \textbf{C}. Large data uses iterative optimizers such as L-BFGS. (A), (B), (D) are wrong.
		
		\item \textbf{E}. Spark bins continuous features for distributed efficiency. (A), (B), (C), (D) are not the reason.
		
		\item \textbf{A}. \texttt{fmin} minimizes the objective. (B), (C), (D) are wrong.
		
		\item \textbf{A}. \texttt{SparkTrials} distributes trials. (B), (C), (D) are wrong.
		
		\item \textbf{C}. \texttt{tpe.suggest} proposes new configurations. (A), (B), (D) are other responsibilities.
		
		\item \textbf{C}. Hyperopt with \texttt{SparkTrials} and TPE leverages the cluster. (A), (B), (D) are inefficient.
		
		\item \textbf{C}. $3 \times 2 \times 3 = 18$. (A), (B), (D) are wrong counts.
		
		\item \textbf{C}. Six combinations can run in parallel. (A), (B), (D) are wrong counts.
		
		\item \textbf{A}. Hyperopt explores more efficiently with TPE. (B), (C), (D) are false.
		
		\item \textbf{A}. StringIndexer converts categories to indices. (B), (C), (D) are other transformers.
		
		\item \textbf{B}. Train-validation trains fewer models. (A), (C), (D) are false.
		
		\item \textbf{B}. Put the pipeline inside the CV when it contains estimators. (A), (C), (D) are not the general rule.
		
		\item \textbf{E}. Recall maximizes identification of positives. (A), (B), (C), (D) measure other things.
		
		\item \textbf{C}. F1 is for imbalance with false-negative priority. (A), (B), (D) are wrong conditions.
		
		\item \textbf{B}. Evaluate precision, recall, F1, and PR-AUC. (A), (C), (D) are wrong actions.
		
		\item \textbf{B}. $R^2$ is the fraction of variance explained. (A), (C), (D) are wrong interpretations.
		
		\item \textbf{A}. RMSE penalizes large errors more. (B), (C), (D) are false.
		
		\item \textbf{B}. Large train-validation gap means overfitting; reduce complexity. (A), (C), (D) are wrong.
		
		\item \textbf{A}. Complexity lowers bias and raises variance. (B), (C), (D) are false.
		
		\item \textbf{E}. RMSE is valid for regression. (A), (B), (C), (D) are possible explanations.
		
		\item \textbf{A}. $(10+12+17)/3 = 13.0$. (B), (C), (D) are wrong.
		
		\item \textbf{B}. Feature engineering on validation data leaks. (A), (C), (D) are not the consequence.
		
		\item \textbf{A}. \texttt{CrossValidator} does k-fold CV with a parameter grid. (B), (C), (D) are other tools.
		
		\item \textbf{B}. A split is cheaper than k-fold CV. (A), (C), (D) are false.
		
		\item \textbf{A}. Feature engineering and dimensionality reduction mitigate overfitting. (B) worsens it; (C) is insufficient; (D) is false.
		
		\item \textbf{A}. Small $k$ means low bias and high variance. (B), (C), (D) are false.
		
		\item \textbf{B}. RBF is common for non-linear SVM. (A) is linear; (C), (D) are not kernels.
		
		\item \textbf{B}. Boosting trains sequentially. (A) is bagging; (C), (D) are other methods.
		
		\item \textbf{A}. \texttt{max\_depth} limits depth. (B), (C), (D) are other hyperparameters.
		
		\item \textbf{A}. Batch, real-time, and streaming are the three serving approaches. (B), (C), (D) are not the exam's three.
		
		\item \textbf{B}. \texttt{spark\_udf} distributes inference. (A), (C), (D) do not scale.
		
		\item \textbf{A}. Streaming inference uses Delta Live Tables. (B), (C), (D) are not streaming inference.
		
		\item \textbf{A}. Traffic splitting is configured on the endpoint. (B), (C), (D) do not split traffic.
		
		\item \textbf{B}. Real-time endpoints serve low-latency REST predictions. (A), (C), (D) are not the purpose.
		
		\item \textbf{D}. Scoring every record with all three models multiplies cost. (A), (B), (C) do not necessarily reduce efficiency.
		
		\item \textbf{A}. Batch inference scores offline and writes predictions. (B), (C), (D) are other concepts.
		
		\item \textbf{B}. A Standard cluster supports SQL and Spark ML. (A) is single-node; (C) is SQL-only; (D) is false.
		
		\item \textbf{A}. Drift degrades deployed models. (B), (C), (D) are false.
		
		\item \textbf{A}. Data drift changes inputs; concept drift changes the relationship. (B), (C), (D) are false.
		
		\item \textbf{A}. Biased data can cause discriminatory outcomes. (B), (C), (D) are not fairness concerns.
		
		\item \textbf{B}. A Pipeline fit on training data prevents leakage. (A), (C), (D) cause leakage.
		
	\end{enumerate}
	
\end{document}